% PAGES: 254-256
% GAP NOTICE: chunk c1's delivered files stop at PDF page 253 (sec08_c1_11.tex,
% folio 237, closing an \answerx block); PDF page 254 (folio 238) was never
% transcribed by c1 and ch08/ contained no file covering it. I (chunk c2,
% assigned 255-266) rendered PDF page 254 directly via CropBox at 300 dpi and
% confirmed it holds the Lemma 8.1 statement and the opening of its proof
% (through eq. (8.61)), which flows directly into what was originally my
% first assigned page, 255 (the "Let b = (b_i)..." paragraph below picks up
% exactly where (8.61) leaves off). I am transcribing page 254 here to close
% this gap -- as the last chunk of the chapter I must leave nothing broken --
% which extends my file range from 255-266 to 254-266. Equation counter
% cross-check: c1's last numbered equation (sec08_c1_10.tex) is (8.56)
% ("The solution (8.56) from section 8.2..."), so the Lemma 8.1 statement
% below correctly auto-numbers to (8.57)-(8.58), matching the printed page.
% No theorem-like environment appears anywhere in c1's files, so this Lemma
% is the first \begin{lemma} in the chapter and auto-numbers to "Lemma 8.1",
% matching print.

Let $\mu_Y$ and $\mu_{X_k}$ be the means of $Y$ and $X_k$, for $k=1:n$, respectively, and
let $\mu_{\bfx}$ be the $n\times1$ column vector $\mu_{\bfx} = (\mu_{X_k})_{k=1:n}$.

\begin{lemma}
Let $X_1, X_2, \ldots, X_n$ be random variables on the same probability space with
nonsingular covariance matrix $\Sigma_{\bfx}$, and let $Y$ be another random variable on
the same probability space. The best approximation of $Y$ by a linear combination of
$X_1, X_2, \ldots, X_n$ plus a constant, i.e., such that the random variable
\[
Y \;-\; \left(a+\sum_{i=1}^n b_iX_i\right) \quad \text{has mean $0$ and minimal
variance,}
\]
where $a, b_1, b_2, \ldots, b_n$ are constants, is given by
\begin{align}
\bfb &= (\Sigma_{\bfx})^{-1}\sigma_{Y,\bfx}.\\
a &= \mu_Y - \mu_{\bfx}^t(\Sigma_{\bfx})^{-1}\sigma_{Y,\bfx},
\end{align}
where $\bfb = (b_i)_{i=1:n}$ and $\sigma_{Y,\bfx} = (\cov(Y,X_i))_{i=1:n}$ is the
covariance vector of $Y$ and $X_1, X_2, \ldots, X_n$.
\end{lemma}

\begin{proof}
Denote by $\mu_Y$ and $\mu_{X_k}$ the mean of $Y$ and of $X_k$, respectively, for
$k=1:n$. Then,
\[
\E{Y - \left(a+\sum_{i=1}^n b_iX_i\right)} \;=\; \mu_Y - a - \sum_{i=1}^n b_i\mu_{X_i},
\]
and therefore $\E{Y-\left(a+\sum_{i=1}^n b_iX_i\right)} = 0$ if and only if
\begin{equation}
a \;=\; \mu_Y - \sum_{i=1}^n b_i\mu_{X_i}.
\end{equation}

Note that
\begin{align*}
\var\left(Y-\left(a+\sum_{i=1}^n b_iX_i\right)\right) &= \var\left(\left(Y-\sum_{i=1}^n
b_iX_i\right)-a\right)\\
&= \var\left(Y-\sum_{i=1}^n b_iX_i\right),
\end{align*}
since $\var(X-a) = \var(X)$ for any random variable $X$ and for any constant $a$.

Thus, we are looking for constants $b_i$, $i=1:n$, such that
\begin{equation}
\var\left(Y-\sum_{i=1}^n b_iX_i\right) \quad \text{is minimal.}
\end{equation}

Recall that, for any random variables $X$ and $Y$ over the same probability space,
\[
\var(Y-X) \;=\; \var(Y) - 2\cov(Y,X) + \var(X).
\]
Then,
\begin{align}
\var\left(Y-\sum_{i=1}^n b_iX_i\right) &= \var(Y) - 2\cov\left(Y,\sum_{i=1}^n
b_iX_i\right) + \var\left(\sum_{i=1}^n b_iX_i\right) \notag\\
&= \var(Y) - 2\sum_{i=1}^n b_i\cov(Y,X_i) + \var\left(\sum_{i=1}^n b_iX_i\right).
\end{align}

Let $\bfb = (b_i)_{i=1:n}$ be the coefficients column vector, and let
\[
\sigma_{Y,\bfx} \;=\; (\cov(Y,X_i))_{i=1:n} \;=\; \begin{pmatrix} \cov(Y,X_1) \\ \vdots \\
\cov(Y,X_n) \end{pmatrix}
\]
be the covariance vector of $Y$ and $X_1, X_2, \ldots, X_n$.

From (1.5), we find that
\begin{equation}
\sum_{i=1}^n b_i\cov(Y,X_i) \;=\; \sigma_{Y,\bfx}^t\bfb.
\end{equation}

Let $\Sigma_{\bfx}$ be the covariance matrix of the random variables $X_1, X_2, \ldots,
X_n$, and recall from (7.25) that
\begin{equation}
\var\left(\sum_{i=1}^n b_iX_i\right) \;=\; \bfb^t\Sigma_{\bfx}\bfb.
\end{equation}

From (8.61--8.63), we obtain that
\begin{equation}
\var\left(Y-\sum_{i=1}^n b_iX_i\right) \;=\; \var(Y) - 2\sigma_{Y,\bfx}^t\bfb +
\bfb^t\Sigma_{\bfx}\bfb.
\end{equation}

Thus, minimizing $\var\left(Y-\sum_{i=1}^n b_iX_i\right)$, see (8.60), is equivalent to
finding a vector $\bfb\in\R^n$ such that
\begin{equation}
\var(Y) - 2\sigma_{Y,\bfx}^t\bfb + \bfb^t\Sigma_{\bfx}\bfb \quad \text{is minimal.}
\end{equation}

Let $f:\R^n\to\R$ be the function given by
\[
f(x) \;=\; \var(Y) - 2\sigma_{Y,\bfx}^tx + x^t\Sigma_{\bfx}x, \quad \forall\, x\in\R^n.
\]
Then, solving (8.65) is equivalent to finding a minimum point for $f(x)$.

Recall from (10.43) and (10.47) the following gradient formulas:
\[
\begin{aligned}
D(C^tx) &= C^t; \\
D\left(x^tAx\right) &= 2(Ax)^t,
\end{aligned}
\]
for any constant column vector $C$ of size $n$ and for any $n\times n$ symmetric matrix
$A$.

Then, the gradient $Df(x)$ of $f(x)$ is
\begin{align}
Df(x) &= D(\var(Y)) - 2D(\sigma_{Y,\bfx}^tx) + D(x^t\Sigma_{\bfx}x) \notag\\
&= -2\sigma_{Y,\bfx}^t + 2(\Sigma_{\bfx}x)^t \notag\\
&= 2(\Sigma_{\bfx}x - \sigma_{Y,\bfx})^t;
\end{align}
note that $D(\var(Y)) = 0$ since $\var(Y)$ is not a function of $x$.

Any minimum point $x_0$ of $f(x)$ is also a critical point of $f(x)$, i.e., a solution to
$Df(x_0) = 0$. From (8.66), we find that, if $Df(x_0) = 0$, then $\Sigma_{\bfx}x_0 =
\sigma_{Y,\bfx}$, which, since we assumed that the covariance matrix $\Sigma_{\bfx}$ is
nonsingular, has the unique solution
\begin{equation}
x_0 \;=\; (\Sigma_{\bfx})^{-1}\sigma_{Y,\bfx}.
\end{equation}

To classify the critical point $x_0$ given by (8.67), we compute the Hessian of $f(x)$.
Recall from (10.45) and (10.48), respectively, that $D^2(C^tx) = 0$ and
$D^2\left(x^tAx\right) = 2A$. Thus, $D^2f(x) = 2\Sigma_{\bfx}$, and therefore $D^2f(x_0)
= 2\Sigma_{\bfx}$.

The covariance matrix $\Sigma_{\bfx}$ is symmetric positive definite since it is
nonsingular. Then, $x_0$ is a local minimum point for the function $f(x)$, and since
$x_0$ is the only critical point of the function $f(x)$, we conclude that $x_0$ is a
global minimum point for the function $f(x)$.

Thus, the vector $\bfb$ minimizing (8.65) is $\bfb = x_0$ given by (8.67), i.e.,
\begin{equation}
\bfb \;=\; (\Sigma_{\bfx})^{-1}\sigma_{Y,\bfx}.
\end{equation}

Let $\mu_{\bfx} = (\mu_{X_i})_{i=1:n}$. Then, from (8.59) and (8.68), we find that
\begin{align}
a &= \mu_Y - \sum_{i=1}^n b_i\mu_i \;=\; \mu_Y - \mu_{\bfx}^t\bfb \notag\\
&= \mu_Y - \mu_{\bfx}^t(\Sigma_{\bfx})^{-1}\sigma_{Y,\bfx}.
\end{align}

From (8.68) and (8.69), we conclude that the constants $a, b_1, b_2, \ldots, b_n$ such
that the random variable $Y - \left(a+\sum_{i=1}^n b_iX_i\right)$ has mean $0$ and
minimal variance are given by (8.57) and (8.58).
\end{proof}

Note that the formulas (8.57) and (8.58) for the coefficients of ordinary least squares
for random variables are the random variable versions of the formulas (8.80) and (8.81)
for time series data.

\section{The intuition behind ordinary least squares for time series data}

The solution (8.56) from section 8.2 for the linear regression of time series data,
while applicable in practice, lacks intuition. We provide this intuition here by
expressing $a$ and $\bfb$ from (8.56) in terms of the time series data for $Y$ and for
$X_k$, $k=1:n$, and by highlighting the connection to the solution to the ordinary least
squares problem for random variables; see Lemma 8.2 and Lemma 8.1.

Let $\widehat\mu_Y$ and $\widehat\mu_{X_k}$ be the sample means of $Y$ and $X_k$, for
$k=1:n$, respectively, i.e.,
\begin{align}
\widehat\mu_Y &= \frac1N\sum_{i=1}^N Y(t_i) \;=\; \frac1N\bfone^tT_Y;\\
\widehat\mu_{X_k} &= \frac1N\sum_{i=1}^N X_k(t_i) \;=\; \frac1N\bfone^tT_{X_k}, \quad
\forall\,k=1:n,
\end{align}
where $T_Y$ and $T_{X_k}$ are the column vectors of the time series data for $Y$ and
$X_k$ given by (8.51) and (8.52), respectively, and let $\widehat\mu_{\bfx}$ be the
$n\times1$ column vector
\begin{equation}
\widehat\mu_{\bfx} \;=\; (\widehat\mu_{X_k})_{k=1:n}.
\end{equation}
